% !TEX root = ../jalase06.tex
%% ============================================================
%%  جلسهٔ ۶ — توزیع هندسی و بی‌حافظگی، توزیع پواسون و ارتباطش با
%%  برنولی و دوجمله‌ای
%% ============================================================
\chapter{توزیع‌های هندسی و پواسون}
\label{ch:s6}

% نمودار میله‌ای PMF:  \stems{رنگ}{k/p, k/p, ...}  (داخل tikzpicture با مقیاس x و y دلخواه)
\providecommand{\stems}[2]{%
  \edef\stemlist{#2}%
  \foreach \k/\p in \stemlist {
    \draw[#1, line width=1.5pt] (\k,0) -- (\k,\p);
    \fill[#1] (\k,\p) circle (1.7pt);
  }}

\sessioncard{دوشنبه ۲۰ مهر}
{توزیع هندسی و خاصیت بی‌حافظگی؛ توزیع پواسون و پیوندش با برنولی و دوجمله‌ای}
{\lr{Pishro-Nik} بخش \lr{3.1.5} \,؛\, \lr{Ross} بخش‌های \lr{4.7} و \lr{4.8.1} \,؛\, \lr{Bertsekas} بخش \lr{2.2}}
{\begin{itemize}
  \item «داستان» توزیع هندسی را بشناسید، \lr{PMF} و \lr{tail} آن، $\P(X>k)=(1-p)^k$، را به دست آورید و در مسئله تشخیصش دهید؛
  \item خاصیت بی‌حافظگی را دقیق بیان و اثبات کنید و نشان دهید هندسی \emph{تنها} توزیع گسستهٔ بی‌حافظه روی $\{1,2,\dots\}$ است؛
  \item توزیع پواسون را به‌عنوان حدِ دوجمله‌ای وقتی $n\to\infty$ و $np=\lambda$ ثابت است به دست آورید و از آن برای تقریب استفاده کنید؛
  \item توضیح دهید چرا جمعِ تعداد زیادی برنولیِ «نادر» تقریباً پواسون است و چه زمانی این تقریب خوب نیست.
\end{itemize}}

در جلسهٔ پیش دو توزیع گسستهٔ مهم، برنولی و دوجمله‌ای، را دیدیم. هر دو از یک منبع می‌آمدند: \key{دنباله‌ای از آزمایش‌های برنولیِ مستقل}. در این جلسه از همان منبع دو توزیع دیگر می‌سازیم: اگر بپرسیم «اولین موفقیت کِی رخ می‌دهد؟» به توزیع \emph{هندسی} می‌رسیم، و اگر تعداد آزمایش‌ها خیلی زیاد و احتمال موفقیت خیلی کم باشد، به توزیع \emph{پواسون}.

%% ------------------------------------------------------------
\section{توزیع هندسی}
\label{sec:geom}
%% ------------------------------------------------------------

\textbf{داستان.} سکه‌ای داریم که با احتمال $p$ شیر می‌آید. آن را \emph{آن‌قدر} پرتاب می‌کنیم تا \emph{اولین شیر} بیاید و $X$ را تعداد کل پرتاب‌ها می‌گیریم. به بیان کلی‌تر: آزمایش‌های برنولیِ مستقل با احتمال موفقیت $p$ را تکرار می‌کنیم تا اولین موفقیت رخ دهد. برد $X$ برابر $\{1,2,3,\dots\}$ است؛ نامتناهی ولی شمارا.

پیشامد $\{X=k\}$ یعنی $k-1$ آزمایشِ اول شکست بخورند و آزمایش $k$-ام موفق شود: $\underbrace{F\,F\cdots F}_{k-1}\,S$. به دلیل استقلال، احتمال این دنبالهٔ مشخص $(1-p)^{k-1}p$ است.

\begin{definition}[label=def:geom]{توزیع هندسی}
متغیر تصادفی $X$ با برد $\{1,2,3,\dots\}$ و
\[ p_X(k)=\P(X=k)=(1-p)^{k-1}\,p,\qquad k=1,2,3,\dots\qquad(0<p\le1) \]
\term{هندسی}{geometric} با پارامتر $p$ نامیده می‌شود و می‌نویسیم $X\sim\mathrm{Geom}(p)$.
\end{definition}

این یک \lr{PMF} معتبر است: احتمال‌ها نامنفی‌اند و با جمع سری هندسی (که نام توزیع از آن آمده است)
\[ \sum_{k=1}^{\infty}(1-p)^{k-1}p=p\sum_{j=0}^{\infty}(1-p)^{j}=p\cdot\frac{1}{1-(1-p)}=1. \]

\begin{figblock}
\begin{tikzpicture}[x=0.42cm, y=5.2cm]
  \foreach \sh/\pp/\col/\data in {
     0/0.3/setA/{1/0.3,2/0.21,3/0.147,4/0.1029,5/0.072,6/0.0504,7/0.0353,8/0.0247,9/0.0173,10/0.0121,11/0.0085,12/0.0059,13/0.0042,14/0.0029},
     18/0.7/accent/{1/0.7,2/0.21,3/0.063,4/0.0189,5/0.0057,6/0.0017,7/0.0005}} {
    \begin{scope}[xshift=\sh*0.42cm]
      \draw[-{Stealth}, graytext] (0,0) -- (15.2,0) node[right, font=\small] {$k$};
      \foreach \tk in {1,5,10} { \draw[graytext] (\tk,0) -- (\tk,-0.015); \node[below, font=\scriptsize] at (\tk,-0.012) {$\tk$}; }
      \stems{\col}{\data}
      \node[font=\small] at (7.5,0.78) {$\mathrm{Geom}(\pp)$};
    \end{scope}
  }
\end{tikzpicture}

\figcaption[fig:geom-pmf]{\lr{PMF} توزیع هندسی. هر میله $(1-p)$ برابرِ میلهٔ قبلی است؛ پس نمودار همیشه نزولی است و قله‌اش، برای \emph{هر} $p$، در $k=1$ قرار دارد. هرچه $p$ کوچک‌تر باشد، احتمال آهسته‌تر کم می‌شود و \lr{tail} توزیع بلندتر است.}
\end{figblock}

\begin{note}
نسبت دو احتمال متوالی همیشه $p_X(k+1)/p_X(k)=1-p<1$ است؛ پس \key{محتمل‌ترین مقدار $X$ همیشه ۱ است}، حتی وقتی $p$ خیلی کوچک باشد. مثلاً با $p=0.01$، احتمال اینکه دقیقاً در پرتاب اول موفق شوید ($0.01$) از احتمال اینکه دقیقاً در پرتاب صدم موفق شوید ($0.99^{99}\cdot0.01\approx0.0037$) بیشتر است؛ هرچند احتمال اینکه دست‌کم ۱۰۰ پرتاب لازم شود حدود $0.37$ است (ببینید پایین‌تر).
\end{note}

\begin{warning}
بعضی کتاب‌ها هندسی را «تعداد \emph{شکست‌ها} قبل از اولین موفقیت» تعریف می‌کنند. با آن تعریف برد $\{0,1,2,\dots\}$ و $\P(Y=k)=(1-p)^kp$ است و $Y=X-1$. در این جزوه (مثل \lr{Pishro-Nik} و \lr{Ross}) همیشه منظور «تعداد کل آزمایش‌ها تا اولین موفقیت» است. مهم این نیست که فرمول را حفظ کنید؛ مهم این است که بتوانید آن را از داستان بسازید.
\end{warning}

\subsection{\lr{tail} توزیع و \lr{CDF}}

برای هندسی، احتمالِ «دیر شدن» از خودِ \lr{PMF} ساده‌تر است:

\begin{proposition}[label=prop:geom-tail]{\lr{tail} توزیع هندسی}
اگر $X\sim\mathrm{Geom}(p)$، برای هر عدد صحیح $k\ge0$:
\[ \P(X>k)=(1-p)^k,\qquad F_X(k)=\P(X\le k)=1-(1-p)^k. \]
\end{proposition}

\begin{proof}
$\{X>k\}$ یعنی «در $k$ آزمایش اول هیچ موفقیتی رخ نداده است»؛ یعنی $k$ شکست پشت سر هم، که به دلیل استقلال احتمالش $(1-p)^k$ است. (همین را می‌توان با جمع سری $\sum_{j>k}(1-p)^{j-1}p$ هم به دست آورد.)
\end{proof}

\begin{example}[label=ex:geom-urn]{برداشتن مهره با جایگذاری (\lr{Ross})}
ظرفی $N$ مهرهٔ سفید و $M$ مهرهٔ سیاه دارد. مهره‌ها را یکی‌یکی و \emph{با جایگذاری} برمی‌داریم تا اولین مهرهٔ سیاه بیاید. احتمال اینکه (الف) دقیقاً $n$ برداشت لازم باشد، و (ب) دست‌کم $k$ برداشت لازم باشد چقدر است؟
\begin{solution}
چون با جایگذاری برمی‌داریم، برداشت‌ها مستقل‌اند و در هر کدام احتمال سیاه بودن $p=\frac{M}{M+N}$ است؛ پس تعداد برداشت‌ها $X\sim\mathrm{Geom}(p)$ است:
\[ \fa{(الف)}\ \ \P(X=n)=\Bigl(\frac{N}{M+N}\Bigr)^{n-1}\frac{M}{M+N},\qquad \fa{(ب)}\ \ \P(X\ge k)=\P(X>k-1)=\Bigl(\frac{N}{M+N}\Bigr)^{k-1}. \]
دقت کنید که «دست‌کم $k$» یعنی $X>k-1$، نه $X>k$.
\end{solution}
\end{example}



%% ------------------------------------------------------------
\section{بی‌حافظگی}
\label{sec:memoryless}
%% ------------------------------------------------------------

فرض کنید تا الان $m$ بار سکه انداخته‌اید و هنوز شیر نیامده است. آیا شیر «طلبکار» شده و احتمال آمدنش بیشتر است؟ شهود می‌گوید نه: سکه حافظه ندارد و پرتاب‌های بعدی از گذشته مستقل‌اند. پس از این لحظه به بعد، انگار همه چیز از نو شروع شده است. به بیان دقیق‌تر: \key{بعد از $n$ پرتاب ناموفق، احتمال‌های پرتاب $(n+1)$-ام به بعد درست مثل پرتاب اول به بعد است}.

\begin{theorem}[label=thm:memoryless]{بی‌حافظگی توزیع هندسی}
اگر $X\sim\mathrm{Geom}(p)$، برای هر دو عدد صحیح نامنفی $m$ و $n$:
\[ \boxed{\ \P(X>m+n\mid X>m)=\P(X>n)\ } \]
و به طور معادل، برای هر $k\ge1$: \ $\P(X=m+k\mid X>m)=\P(X=k)=(1-p)^{k-1}p$.
\end{theorem}

\begin{proof}
چون $\{X>m+n\}\subseteq\{X>m\}$، اشتراک این دو پیشامد همان $\{X>m+n\}$ است. با گزارهٔ \ref{prop:geom-tail}:
\[ \P(X>m+n\mid X>m)=\frac{\P(X>m+n)}{\P(X>m)}=\frac{(1-p)^{m+n}}{(1-p)^m}=(1-p)^n=\P(X>n). \]
برای صورت دوم: $\{X=m+k\}\subseteq\{X>m\}$ و
$\P(X=m+k\mid X>m)=\frac{(1-p)^{m+k-1}p}{(1-p)^m}=(1-p)^{k-1}p$.
\end{proof}

\begin{intuition}
صورت دوم قضیه را این‌طور بخوانید: \emph{به شرط اینکه $m$ پرتاب اول ناموفق بوده باشد، تعداد پرتاب‌های \emph{باقی‌مانده}، یعنی $X-m$، باز هم $\mathrm{Geom}(p)$ است}. اگر کسی بعد از $m$ پرتاب ناموفق وارد اتاق شود و شروع به شمردن کند، چیزی که می‌بیند از نظر احتمالاتی با کسی که از اول آنجا بوده فرقی ندارد. توزیعِ «زمان انتظارِ باقی‌مانده» به اینکه چقدر صبر کرده‌ایم بستگی ندارد.
\end{intuition}

\begin{figblock}
\begin{tikzpicture}[x=0.5cm, y=4.2cm]
  \draw[-{Stealth}, graytext] (0,0) -- (17,0) node[right, font=\small] {$k$};
  \fill[graytext!12] (0.5,0) rectangle (4.5,0.36);
  \node[font=\scriptsize, graytext] at (2.5,0.39) {\rl{$m=4$ شکست}};
  \stems{graytext!50}{1/0.3,2/0.21,3/0.147,4/0.1029}
  \begin{scope}[xshift=-2.2pt]\stems{setA}{5/0.072,6/0.0504,7/0.0353,8/0.0247,9/0.0173,10/0.0121,11/0.0085,12/0.0059,13/0.0042,14/0.0029,15/0.002}\end{scope}
  \begin{scope}[xshift=2.2pt]\stems{accent}{5/0.3,6/0.21,7/0.147,8/0.1029,9/0.072,10/0.0504,11/0.0353,12/0.0247,13/0.0173,14/0.0121,15/0.0085}\end{scope}
  \foreach \tk in {1,5,10,15} \node[below, font=\scriptsize] at (\tk,0) {$\tk$};
  \node[font=\small, setA!80!black, anchor=south west] at (10.6,0.035) {$\P(X=k)$};
  \node[font=\small, accent!80!black, anchor=west] at (7.2,0.3) {$\P(X=k\mid X>4)$};
\end{tikzpicture}

\figcaption[fig:memoryless]{بی‌حافظگی برای $\mathrm{Geom}(0.3)$: اگر بدانیم $X>4$، احتمال‌های باقی‌مانده (آبی) بر $\P(X>4)$ تقسیم می‌شوند و نمودار نارنجی به دست می‌آید که دقیقاً همان نمودار اصلی است، فقط ۴ واحد به راست رفته است.}
\end{figblock}

\begin{example}[label=ex:roulette]{مغالطهٔ قمارباز}
در رولت اروپایی ۳۷ خانه هست که ۱۸ تای آن قرمز است. قماربازی ۱۰ دور متوالی دیده که قرمز نیامده و فکر می‌کند «حالا دیگر نوبت قرمز است». احتمال اینکه در ۳ دور بعدی هم قرمز نیاید چقدر است؟ اگر کسی تازه سر میز آمده باشد چطور؟
\begin{solution}
تعداد دورها تا اولین قرمز $X\sim\mathrm{Geom}(\tfrac{18}{37})$ است. طبق بی‌حافظگی
\[ \P(X>13\mid X>10)=\P(X>3)=\Bigl(\frac{19}{37}\Bigr)^3\approx0.135, \]
\emph{دقیقاً همان} احتمالی که تازه‌وارد دارد. ۱۰ دور قرمز نیامدن هیچ چیزی را «جبران» نمی‌کند.
\end{solution}
\end{example}


\begin{example}[label=ex:keys]{کلیدها: با حافظه و بی‌حافظه}
کسی $n$ کلید دارد که فقط یکی از آن‌ها در را باز می‌کند. $X$ را تعداد تلاش‌ها تا باز شدن در بگیرید. (الف) اگر هر بار یک کلید را کاملاً تصادفی از بین هر $n$ کلید امتحان کند (کلیدهای امتحان‌شده را کنار نگذارد)، و (ب) اگر کلیدهای امتحان‌شده را کنار بگذارد، \lr{PMF} متغیر $X$ را بیابید. کدام بی‌حافظه است؟
\begin{solution}
(الف) تلاش‌ها مستقل‌اند و هر کدام با احتمال $\tfrac1n$ موفق است؛ پس $X\sim\mathrm{Geom}(\tfrac1n)$ و $X$ بی‌حافظه است.\\
(ب) حالا تلاش‌ها مستقل نیستند. با قاعدهٔ ضرب (جلسهٔ ۴) برای $k=1,\dots,n$:
\[ \P(X=k)=\frac{n-1}{n}\cdot\frac{n-2}{n-1}\cdots\frac{n-k+1}{n-k+2}\cdot\frac{1}{n-k+1}=\frac1n. \]
پس هر یک از مقادیر $1,\dots,n$ هم‌شانس است و $\P(X>k)=\frac{n-k}{n}$. این توزیع حافظه دارد:
\[ \P(X>k+1\mid X>k)=\frac{n-k-1}{n-k}, \]
که با افزایش $k$ \emph{کم} می‌شود: هر شکست، شانس موفقیت در تلاش بعدی را بیشتر می‌کند، و بعد از $n-1$ شکست، تلاش بعدی حتماً موفق است. در حالت (الف) برعکس، با احتمال $(1-\frac1n)^n\approx e^{-1}\approx0.37$ حتی $n$ تلاش هم کافی نیست!
\end{solution}
\end{example}

حالا فرمول دقیقِ بی‌حافظگی را به عکس به کار می‌بریم و نشان می‌دهیم این خاصیت، توزیع هندسی را \emph{کاملاً} مشخص می‌کند:

\begin{theorem}[label=thm:memoryless-unique]{هندسی تنها توزیع گسستهٔ بی‌حافظه است}
فرض کنید $X$ متغیری با برد $\{1,2,3,\dots\}$ باشد، $\P(X>n)>0$ برای هر $n$، و برای هر دو عدد صحیح نامنفی $m,n$:
\[ \P(X>m+n\mid X>m)=\P(X>n). \]
آنگاه $X\sim\mathrm{Geom}(p)$ با $p=\P(X=1)$.
\end{theorem}

\begin{proof}
$G(n)\defeq\P(X>n)$ را تعریف کنید؛ $G(0)=1$. شرط بی‌حافظگی یعنی $G(m+n)/G(m)=G(n)$، یا
\[ G(m+n)=G(m)\,G(n). \]
با استقرا $G(n)=G(1)^n$. قرار دهید $q=G(1)=1-\P(X=1)$. طبق فرض $q>0$؛ و $q<1$، چون اگر $q=1$ بود، $G(n)=1$ برای همهٔ $n$ و این با $\lim_nG(n)=0$ (پیوستگی احتمال) تناقض داشت. حالا
\[ \P(X=k)=G(k-1)-G(k)=q^{k-1}-q^k=q^{k-1}(1-q), \]
که همان \lr{PMF} هندسی با $p=1-q$ است.
\end{proof}

\begin{note}
معادلهٔ $G(m+n)=G(m)G(n)$ می‌گوید «احتمال زنده ماندن» باید \emph{نمایی} از زمان باشد. در جلسهٔ بعد خواهیم دید که در دنیای پیوسته همین معادله به $\P(X>t)=e^{-\lambda t}$ می‌رسد و توزیع \emph{نمایی} را می‌سازد؛ به همین دلیل می‌توان نمایی را «هندسیِ پیوسته‌شده» دانست.
\end{note}

%% ------------------------------------------------------------
\section{توزیع پواسون}
\label{sec:poisson}
%% ------------------------------------------------------------

\subsection{انگیزه: درخواست‌های گوگل}

فرض کنید $n=10^{10}$ دستگاه در دنیا هست و هر کدام، مستقل از بقیه، در یک ثانیهٔ مشخص با احتمال $p=10^{-5}$ یک درخواست به سرورهای گوگل می‌فرستد. تعداد درخواست‌ها در این ثانیه
\[ X\sim\mathrm{Bin}(10^{10},\,10^{-5}) \]
است. اما حساب کردن $\binom{10^{10}}{k}p^k(1-p)^{10^{10}-k}$ عملاً ناممکن است. خوشبختانه در این وضعیت، یعنی «$n$ خیلی بزرگ و $p$ خیلی کوچک»، فرمول خیلی ساده می‌شود. اگر $k$ در مقایسه با $n$ کوچک باشد:
\[ \binom nk=\frac{n(n-1)\cdots(n-k+1)}{k!}\approx\frac{n^k}{k!},\qquad (1-p)^{n-k}\approx(1-p)^n\approx e^{-np}, \]
که دومی از $\ln(1-p)\approx-p$ برای $p$ کوچک می‌آید. پس با $\lambda\defeq np$:
\[ \P(X=k)\approx\frac{n^k}{k!}\,p^k\,e^{-np}=e^{-\lambda}\frac{\lambda^k}{k!}. \]
نکتهٔ جالب این است که جواب فقط به \emph{حاصل‌ضرب} $\lambda=np$ وابسته است، نه به $n$ و $p$ جداگانه.

\subsection{تعریف و حدِ دوجمله‌ای}

\begin{definition}[label=def:poisson]{توزیع پواسون}
متغیر تصادفی $X$ با برد $\{0,1,2,\dots\}$ و
\[ p_X(k)=\P(X=k)=e^{-\lambda}\,\frac{\lambda^k}{k!},\qquad k=0,1,2,\dots \]
\term{پواسون}{Poisson} با پارامتر $\lambda>0$ نامیده می‌شود و می‌نویسیم $X\sim\mathrm{Poisson}(\lambda)$.
\end{definition}

این \lr{PMF} معتبر است: احتمال‌ها مثبت‌اند و با سری تیلور $e^\lambda=\sum_{k\ge0}\frac{\lambda^k}{k!}$:
\[ \sum_{k=0}^{\infty}e^{-\lambda}\frac{\lambda^k}{k!}=e^{-\lambda}\,e^{\lambda}=1. \]

\begin{theorem}[label=thm:poisson-limit]{پواسون، حدِ دوجمله‌ای}
فرض کنید $\lambda>0$ ثابت باشد و $X_n\sim\mathrm{Bin}(n,\frac{\lambda}{n})$. آنگاه برای هر $k\in\{0,1,2,\dots\}$:
\[ \lim_{n\to\infty}\P(X_n=k)=e^{-\lambda}\,\frac{\lambda^k}{k!}. \]
\end{theorem}

\begin{proof}
با $p=\lambda/n$:
\[ \P(X_n=k)=\binom nk\Bigl(\frac\lambda n\Bigr)^k\Bigl(1-\frac\lambda n\Bigr)^{n-k}
=\frac{\lambda^k}{k!}\cdot\underbrace{\frac{n(n-1)\cdots(n-k+1)}{n^k}}_{\to\,1}\cdot\underbrace{\Bigl(1-\frac\lambda n\Bigr)^{n}}_{\to\,e^{-\lambda}}\cdot\underbrace{\Bigl(1-\frac\lambda n\Bigr)^{-k}}_{\to\,1}. \]
عامل اول حاصل‌ضرب $k$ کسرِ $\frac{n-j}{n}$ است که $k$ ثابت است و هر کدام به ۱ میل می‌کند؛ عامل دوم حد معروف ${(1+\frac xn)^n\to e^x}$ است؛ و عامل سوم توانِ ثابتی از عددی است که به ۱ میل می‌کند.
\end{proof}

\begin{keyformula}
\[ \mathrm{Bin}(n,p)\ \approx\ \mathrm{Poisson}(np)\qquad\text{\rl{وقتی $n$ بزرگ، $p$ کوچک و $np$ «متوسط» باشد}} \]
\end{keyformula}

\subsection{شکل نمودار}

\begin{proposition}[label=prop:poisson-mode]{شکل نمودار پواسون}
اگر $X\sim\mathrm{Poisson}(\lambda)$، برای $k\ge1$:
\[ \frac{p_X(k)}{p_X(k-1)}=\frac{\lambda}{k}. \]
پس $p_X(k)$ تا وقتی $k<\lambda$ صعودی و بعد از آن نزولی است؛ قله در $\lfloor\lambda\rfloor$ است (و اگر $\lambda$ صحیح باشد، $\lambda-1$ و $\lambda$ هر دو قله‌اند). همچنین $p_X(0)=e^{-\lambda}$.
\end{proposition}

\begin{proof}
$\frac{e^{-\lambda}\lambda^k/k!}{e^{-\lambda}\lambda^{k-1}/(k-1)!}=\frac\lambda k$، و این نسبت بزرگ‌تر از ۱ است اگر و تنها اگر $k<\lambda$.
\end{proof}

\begin{figblock}
\begin{tikzpicture}[x=0.21cm, y=7cm]
  \foreach \sh/\ll/\col/\data in {
     0/1/setA/{0/0.3679,1/0.3679,2/0.1839,3/0.0613,4/0.0153,5/0.0031,6/0.0005},
     25/4/excolor/{0/0.0183,1/0.0733,2/0.1465,3/0.1954,4/0.1954,5/0.1563,6/0.1042,7/0.0595,8/0.0298,9/0.0132,10/0.0053,11/0.0019,12/0.0006},
     50/10/accent/{1/0.0005,2/0.0023,3/0.0076,4/0.0189,5/0.0378,6/0.0631,7/0.0901,8/0.1126,9/0.1251,10/0.1251,11/0.1137,12/0.0948,13/0.0729,14/0.0521,15/0.0347,16/0.0217,17/0.0128,18/0.0071,19/0.0037,20/0.0019}} {
    \begin{scope}[xshift=\sh*0.21cm]
      \draw[graytext] (-0.5,0) -- (20.8,0);
      \foreach \tk in {0,10,20} { \draw[graytext] (\tk,0) -- (\tk,-0.01); \node[below, font=\scriptsize] at (\tk,-0.008) {$\tk$}; }
      \stems{\col}{\data}
      \node[font=\small] at (10,0.43) {$\mathrm{Poisson}(\ll)$};
    \end{scope}
  }
\end{tikzpicture}

\figcaption[fig:poisson-shapes]{\lr{PMF} پواسون برای $\lambda=1,4,10$. قله نزدیک $\lambda$ است و با بزرگ شدن $\lambda$ نمودار به راست می‌رود، پهن‌تر می‌شود و شکلی زنگوله‌ای پیدا می‌کند؛ درست مثل دوجمله‌ای.}
\end{figblock}

\subsection{داستان پواسون: شمردن پیشامدهای نادر}

قضیهٔ \ref{thm:poisson-limit} داستان پواسون را می‌گوید: \key{تعداد دفعاتی که یک پیشامد نادر رخ می‌دهد، وقتی فرصت‌های بسیار زیاد و تقریباً مستقلی برای رخ دادنش وجود دارد}. نمونه‌ها: تعداد ایمیل‌هایی که در یک دقیقه می‌رسد، تعداد غلط‌های تایپی یک صفحه، تعداد تماس‌ها با یک مرکز تلفن در یک ساعت، تعداد ذرات واپاشیده از یک مادهٔ رادیواکتیو در یک ثانیه، تعداد درخواست‌ها به یک سرور، یا تعداد تصادف‌ها در یک تقاطع در یک ماه. در همهٔ این‌ها تعداد زیادی «فرصت» هست (کاربران، حروف، اتم‌ها، رانندگان) که هر کدام با احتمال کمی پیشامد را می‌سازند.

\begin{note}
\textbf{پارامتر $\lambda$ با طول بازه متناسب است.} اگر پیشامدها به طور متوسط با «نرخ» $\alpha$ در واحد زمان رخ دهند، تعداد آن‌ها در بازه‌ای به طول $t$ تقریباً $\mathrm{Poisson}(\alpha t)$ است: بازهٔ دو برابر، دو برابر فرصت دارد. (در جلسه‌های بعد خواهیم دید که $\lambda$ دقیقاً «میانگینِ» تعداد پیشامدهاست.)
\end{note}

\begin{example}[label=ex:emails]{ایمیل‌ها (\lr{Pishro-Nik})}
تعداد ایمیل‌هایی که در یک روز کاری می‌رسد را با پواسون مدل می‌کنیم؛ به طور متوسط ۰٫۲ ایمیل در دقیقه. (الف) احتمال اینکه در یک بازهٔ ۵ دقیقه‌ای هیچ ایمیلی نرسد چقدر است؟ (ب) احتمال اینکه در ۱۰ دقیقه بیش از ۳ ایمیل برسد؟
\begin{solution}
(الف) تعداد ایمیل‌ها در ۵ دقیقه $X\sim\mathrm{Poisson}(5\times0.2)=\mathrm{Poisson}(1)$ است و $\P(X=0)=e^{-1}\approx0.368$.\\
(ب) در ۱۰ دقیقه $Y\sim\mathrm{Poisson}(2)$ و
\[ \P(Y>3)=1-e^{-2}\Bigl(1+2+\frac{2^2}{2!}+\frac{2^3}{3!}\Bigr)=1-\frac{19}{3e^2}\approx0.143. \]
\end{solution}
\end{example}

\begin{example}[label=ex:server]{ظرفیت سرور}
یک سرور به طور متوسط ۴ درخواست در ثانیه دریافت می‌کند و تعداد درخواست‌ها در هر ثانیه را $\mathrm{Poisson}(4)$ مدل می‌کنیم. سرور باید چند درخواست در ثانیه را بتواند پاسخ دهد تا دست‌کم در ۹۹٪ ثانیه‌ها از کار نیفتد؟
\begin{solution}
کمترین $c$ را می‌خواهیم که $F_X(c)=\P(X\le c)\ge0.99$. \lr{CDF} را پله به پله جمع می‌زنیم:
\begin{center}
\renewcommand{\arraystretch}{1.2}
\small
\begin{tabular}{c|ccccccccccc}
$k$ & $0$ & $1$ & $2$ & $3$ & $4$ & $5$ & $6$ & $7$ & $8$ & $9$ & $10$\\\hline
$p_X(k)$ & $.018$ & $.073$ & $.147$ & $.195$ & $.195$ & $.156$ & $.104$ & $.060$ & $.030$ & $.013$ & $.005$\\
$F_X(k)$ & $.018$ & $.092$ & $.238$ & $.433$ & $.629$ & $.785$ & $.889$ & $.949$ & $.979$ & $\mathbf{.992}$ & $.997$
\end{tabular}
\end{center}
پس $c=9$: ظرفیتی بیش از \emph{دو برابرِ} میانگین لازم است. نکتهٔ مهم این است که سرور را بر اساس میانگین طراحی نمی‌کنند، بلکه طوری که تا جای ممکن \emph{\lr{tail}} توزیع را هم پوشش دهد.
\end{solution}
\end{example}

\begin{example}[label=ex:typos]{غلط تایپی (\lr{Ross})}
تعداد غلط‌های تایپی هر صفحهٔ یک کتاب $\mathrm{Poisson}(\tfrac12)$ است. احتمال اینکه یک صفحه دست‌کم یک غلط داشته باشد چقدر است؟
\begin{solution}
$\P(X\ge1)=1-\P(X=0)=1-e^{-1/2}\approx0.393$.
\end{solution}
\end{example}

%% ------------------------------------------------------------
\section{پیوند پواسون با برنولی و دوجمله‌ای}
\label{sec:poisson-bern}
%% ------------------------------------------------------------

قضیهٔ \ref{thm:poisson-limit} را با زبان جلسهٔ ۵ بازنویسی کنیم. اگر $I_1,\dots,I_n$ نشانگرهای $n$ پیشامد مستقل با احتمال $\frac\lambda n$ باشند، $I_1+\cdots+I_n\sim\mathrm{Bin}(n,\frac\lambda n)$ است. پس:

\begin{intuition}
\textbf{پواسون = جمعِ تعداد زیادی برنولیِ نادر.} اگر دوجمله‌ای مولکولی بود که از $n$ اتم برنولی ساخته شده، پواسون حالتی است که تعداد اتم‌ها بی‌نهایت می‌شود و هر اتم تقریباً همیشه «خاموش» است، ولی تعداد متوسط اتم‌های «روشن»، $np=\lambda$، ثابت می‌ماند. در جدول زیر:
\begin{center}
\renewcommand{\arraystretch}{1.3}
\begin{tabular}{c|c|c}
& \rl{تعداد آزمایش‌ها} & \rl{احتمال موفقیت در هر آزمایش}\\\hline
$\mathrm{Bern}(p)$ & $1$ & $p$\\
$\mathrm{Bin}(n,p)$ & $n$ & $p$\\
$\mathrm{Poisson}(\lambda)$ & $n\to\infty$ & $p=\lambda/n\to0$
\end{tabular}
\end{center}
\end{intuition}

\subsection{از فرایند برنولی به پواسون}

دنباله‌ای نامتناهی از آزمایش‌های برنولیِ مستقل با احتمال $p$ را \term{فرایند برنولی}{Bernoulli process} می‌نامیم. سه توزیعی که تا اینجا دیدیم، سه سؤال متفاوت دربارهٔ \emph{همین} دنباله‌اند:
\begin{itemize}
  \item هر آزمایش به‌تنهایی: برنولی؛
  \item تعداد موفقیت‌ها در $n$ آزمایش اول: دوجمله‌ای؛
  \item تعداد آزمایش‌ها تا اولین موفقیت (و به دلیل بی‌حافظگی، فاصلهٔ بین هر دو موفقیت متوالی): هندسی.
\end{itemize}

\begin{figblock}
\begin{tikzpicture}[x=0.78cm]
  \foreach \i in {1,...,16} {
    \pgfmathtruncatemacro{\s}{(\i==3)||(\i==7)||(\i==8)||(\i==13)}
    \ifnum\s=1
      \node[draw=accent, fill=accent!25, thick, minimum size=6mm, inner sep=0pt, rounded corners=2pt] (b\i) at (\i,0) {\small$S$};
    \else
      \node[draw=graytext, fill=white, thick, minimum size=6mm, inner sep=0pt, rounded corners=2pt] (b\i) at (\i,0) {\small$F$};
    \fi
    \node[font=\tiny, graytext, above=0.5mm of b\i] {\i};
  }
  \node at (17.2,0) {$\cdots$};
  % فاصله‌ها: هندسی
  \foreach \a/\b/\lab in {0.6/3.4/{T_1=3}, 3.6/7.4/{T_2=4}, 7.6/8.4/{T_3=1}, 8.6/13.4/{T_4=5}}
    \draw[decorate, decoration={brace, mirror, amplitude=4pt}, thick, setA] (\a,-0.45) -- (\b,-0.45) node[midway, below=5pt, font=\scriptsize] {$\lab$};
  % شمارش: دوجمله‌ای
  \draw[decorate, decoration={brace, amplitude=5pt}, thick, excolor] (0.6,0.75) -- (10.4,0.75)
     node[midway, above=6pt, font=\small] {\rl{تعداد موفقیت‌ها در ۱۰ آزمایش اول: $3$؛ توزیع این تعداد $\mathrm{Bin}(10,p)$ است}};
  \node[font=\small, setA!80!black] at (8.5,-1.45) {\rl{فاصله‌های بین موفقیت‌ها: هر کدام $\mathrm{Geom}(p)$}};
\end{tikzpicture}

\figcaption[fig:bern-process]{یک برآمد از فرایند برنولی. شمردن موفقیت‌ها در یک پنجرهٔ ثابت، دوجمله‌ای می‌دهد و فاصلهٔ بین موفقیت‌ها هندسی است. (برعکس، با کنار هم گذاشتن فاصله‌های هندسی هم می‌توان نمونه‌ای از دوجمله‌ای ساخت.)}
\end{figblock}

حالا بازهٔ زمانی $[0,t]$ را در نظر بگیرید و آن را به $n$ تکهٔ کوچک به طول $t/n$ تقسیم کنید. فرض کنید در هر تکه، مستقل از بقیه، با احتمال $\alpha t/n$ یک پیشامد (مثلاً رسیدن یک ایمیل) رخ دهد. تعداد پیشامدها در $[0,t]$ برابر $\mathrm{Bin}(n,\alpha t/n)$ است و وقتی تکه‌ها را ریزتر و ریزتر کنیم ($n\to\infty$)، طبق قضیهٔ \ref{thm:poisson-limit} به $\mathrm{Poisson}(\alpha t)$ می‌رسیم. پواسون، نسخهٔ «زمان پیوستهٔ» فرایند برنولی است.

\begin{figblock}
\begin{tikzpicture}[x=0.48cm]
  \draw[thick, primary, -{Stealth}] (0,0) -- (25.6,0) node[right, font=\small] {\rl{زمان}};
  \foreach \i in {0,...,24} \draw[graytext] (\i,-0.12) -- (\i,0.12);
  \foreach \x in {2.5,9.5,10.5,17.5} \fill[accent] (\x,0.32) circle (3.2pt);
  \node[below, font=\small] at (0,-0.15) {$0$};
  \node[below, font=\small] at (24,-0.15) {$t$};
  \draw[decorate, decoration={brace, mirror, amplitude=3pt}, graytext] (5,-0.55) -- (6,-0.55) node[midway, below=4pt, font=\scriptsize] {$t/n$};
  \node[font=\small] at (12,1.45) {\rl{هر تکه یک آزمایش برنولی با احتمال $p=\alpha t/n$ است؛ پس تعداد کل پیشامدها در $[0,t]$:}};
  \node[font=\small] at (12,0.85) {$N_{[0,t]}\ \sim\ \mathrm{Bin}(n,\alpha t/n)\ \xrightarrow{\ n\to\infty\ }\ \mathrm{Poisson}(\alpha t)$};
\end{tikzpicture}

\figcaption[fig:slots]{تقسیم زمان به تکه‌های ریز. نقطه‌های نارنجی تکه‌هایی هستند که در آن‌ها پیشامد رخ داده است.}
\end{figblock}

\subsection{مثال: خطای \lr{5XX}}

\begin{example}[label=ex:5xx]{خطای \lr{5XX}}
در طول شبانه‌روز تعداد درخواست‌های ارسالی به سرور سامانهٔ درس‌افزار از توزیع پواسون با پارامتر $\lambda$ پیروی می‌کند. با توجه به قدرت بالای پاسخ‌دهی سرور(!) هر درخواست به احتمال $p$ پاسخ داده می‌شود و در غیر این صورت یک خطای دلخواه \lr{5XX} ارسال می‌شود (پاسخ‌دهی به درخواست‌های مختلف مستقل از هم است). توزیع متغیر تصادفی $S$، تعداد درخواست‌هایی که به درستی پاسخ خود را گرفته‌اند، را بیابید. مسئله را از طریق تقریب دوجمله‌ای و پواسون حل کنید.
\begin{solution}
\textbf{ایده.} تعداد کل درخواست‌ها $N\sim\mathrm{Poisson}(\lambda)$ است. طبق بخش قبل، پواسون یعنی «شبانه‌روز را به تعداد زیادی تکهٔ ریز تقسیم کن و در هر تکه یک آزمایش برنولیِ نادر انجام بده». پس کافی است ببینیم در هر تکه، پیشامدِ «درخواستی برسد \emph{و} درست پاسخ بگیرد» با چه احتمالی رخ می‌دهد.

\textbf{گام ۱: تقریب دوجمله‌ای برای $N$.} شبانه‌روز را به $n$ تکهٔ مساوی تقسیم کنید ($n$ خیلی بزرگ، به طوری که در هر تکه حداکثر یک درخواست برسد). طبق قضیهٔ \ref{thm:poisson-limit}، $N\sim\mathrm{Poisson}(\lambda)$ همان حدِ $\mathrm{Bin}(n,\frac\lambda n)$ است؛ یعنی می‌توانیم فرض کنیم در تکهٔ $i$، مستقل از بقیه، با احتمال $\frac\lambda n$ یک درخواست می‌رسد. این پیشامد را $A_i$ بنامید.

\textbf{گام ۲: هر تکه یک برنولی جدید.} $B_i$ را پیشامد «درخواستِ تکهٔ $i$ درست پاسخ بگیرد» بگیرید. پاسخ‌دهی سرور از رسیدن درخواست‌ها مستقل است؛ پس نشانگرِ «در تکهٔ $i$ یک درخواستِ درست‌پاسخ‌داده‌شده داریم» برابر است با
\[ J_i=\ind_{A_i\cap B_i}=\ind_{A_i}\,\ind_{B_i}\ \sim\ \mathrm{Bern}\Bigl(\frac{\lambda}{n}\cdot p\Bigr)=\mathrm{Bern}\Bigl(\frac{\lambda p}{n}\Bigr), \]
و چون تکه‌ها مستقل‌اند، $J_1,\dots,J_n$ هم مستقل‌اند.

\textbf{گام ۳: دوجمله‌ای و سپس پواسون.} تعداد درخواست‌های درست‌پاسخ‌داده‌شده جمع این نشانگرهاست؛ طبق قضیهٔ «جمع برنولی‌ها» در جلسهٔ ۵:
\[ S=J_1+J_2+\cdots+J_n\ \sim\ \mathrm{Bin}\Bigl(n,\frac{\lambda p}{n}\Bigr). \]
با $n\to\infty$ و دوباره طبق قضیهٔ \ref{thm:poisson-limit} (این بار با پارامتر $\lambda p$ به جای $\lambda$):
\[ \boxed{\ S\sim\mathrm{Poisson}(\lambda p),\qquad \P(S=k)=e^{-\lambda p}\,\frac{(\lambda p)^k}{k!},\quad k=0,1,2,\dots\ } \]
به همین ترتیب، تعداد خطاهای \lr{5XX} توزیع $\mathrm{Poisson}\bigl(\lambda(1-p)\bigr)$ دارد.
\medskip

\textbf{بررسی با محاسبهٔ دقیق.} همین جواب را بدون تقریب و با قانون احتمال کل (جلسهٔ ۴) روی مقدار $N$ هم می‌توان گرفت. به شرط $N=m$، هر یک از $m$ درخواست مستقلاً با احتمال $p$ درست پاسخ می‌گیرد؛ پس
\[ \P(S=k\mid N=m)=\binom mkp^k(1-p)^{m-k},\qquad m\ge k. \]
در نتیجه
\begin{align*}
\P(S=k)&=\sum_{m=k}^{\infty}e^{-\lambda}\frac{\lambda^m}{m!}\binom mkp^k(1-p)^{m-k}
=e^{-\lambda}\frac{(\lambda p)^k}{k!}\sum_{m=k}^{\infty}\frac{\bigl(\lambda(1-p)\bigr)^{m-k}}{(m-k)!}\\
&=e^{-\lambda}\frac{(\lambda p)^k}{k!}\,e^{\lambda(1-p)}=e^{-\lambda p}\frac{(\lambda p)^k}{k!}.
\end{align*}
(در تساوی دوم از $\binom mk\frac{1}{m!}=\frac{1}{k!\,(m-k)!}$ و در تساوی سوم از سری تیلور $e^x$ استفاده کردیم.)

\textbf{یک عدد.} فرض کنید روزانه به طور متوسط $\lambda=2000$ درخواست برسد و $p=0.999$. تعداد خطاهای \lr{5XX} در یک روز $\mathrm{Poisson}(2000\times0.001)=\mathrm{Poisson}(2)$ است؛ پس احتمال اینکه یک روز کاملاً بدون خطا بگذرد فقط $e^{-2}\approx0.135$ است و احتمال اینکه حداکثر ۲ خطا رخ دهد $e^{-2}(1+2+2)\approx0.677$.
\end{solution}
\end{example}

\begin{intuition}
به این عمل \term{رقیق‌سازی}{thinning} پواسون می‌گویند: اگر هر پیشامدِ یک شمارش پواسون را مستقلاً با احتمال $p$ «نگه داریم»، شمارشِ پیشامدهای نگه‌داشته‌شده باز هم پواسون است و پارامترش $p$ برابر می‌شود. از دید «جمع برنولی‌های نادر» این کاملاً طبیعی است: رقیق‌سازی فقط احتمالِ هر برنولیِ نادر را از $\frac\lambda n$ به $\frac{\lambda p}{n}$ کم می‌کند، و جمعِ برنولی‌های نادر همچنان پواسون است.
\end{intuition}

\subsection{جمع‌بندی توزیع‌ها}

\begin{center}
\renewcommand{\arraystretch}{1.6}
\small
\begin{tabular}{l|l|l|p{4.4cm}}
\textbf{توزیع} & \textbf{برد} & \textbf{\lr{PMF}} & \textbf{داستان}\\\hline
$\mathrm{Geom}(p)$ & $1,2,\dots$ & $(1-p)^{k-1}p$ & \rl{تعداد آزمایش‌ها تا اولین موفقیت؛ بی‌حافظه}\\
$\mathrm{Poisson}(\lambda)$ & $0,1,2,\dots$ & $e^{-\lambda}\frac{\lambda^k}{k!}$ & \rl{تعداد پیشامدهای نادر؛ حد $\mathrm{Bin}(n,\frac\lambda n)$}
\end{tabular}
\end{center}

%% ------------------------------------------------------------
\section*{جمع‌بندی جلسهٔ ششم}
\addcontentsline{toc}{section}{خلاصه و تمرین‌ها}
%% ------------------------------------------------------------

\begin{summary}
\begin{itemize}
  \item \textbf{هندسی:} تعداد آزمایش‌های برنولی مستقل تا اولین موفقیت؛ $p_X(k)=(1-p)^{k-1}p$ و $\P(X>k)=(1-p)^k$. نمودار همیشه نزولی است و قله در ۱ است.
  \item \textbf{بی‌حافظگی:} $\P(X>m+n\mid X>m)=\P(X>n)$؛ به شرط $m$ شکست، تعداد آزمایش‌های باقی‌مانده باز هم $\mathrm{Geom}(p)$ است. هندسی \emph{تنها} توزیع بی‌حافظه روی $\{1,2,\dots\}$ است، چون $G(m+n)=G(m)G(n)$ نتیجه می‌دهد $G(n)=q^n$.
  \item \textbf{پواسون:} $p_X(k)=e^{-\lambda}\lambda^k/k!$ برای $k=0,1,\dots$؛ حد $\mathrm{Bin}(n,\lambda/n)$ وقتی $n\to\infty$. قله نزدیک $\lambda$ است و $\lambda$ با طول بازهٔ زمانی متناسب است.
  \item \textbf{پواسون و برنولی:} جمعِ تعداد زیادی نشانگرِ نادر و تقریباً مستقل تقریباً $\mathrm{Poisson}(\sum p_i)$ است (حتی با احتمال‌های متفاوت، مثلاً مسئلهٔ کلاه‌ها). اگر $p$ کوچک نباشد یا وابستگی قوی باشد، تقریب بد است.
\end{itemize}
\end{summary}

\subsection*{تمرین‌های پیشنهادی}

\begin{exercise}
اگر $X\sim\mathrm{Geom}(p)$، احتمال اینکه $X$ زوج باشد چقدر است؟ آیا این احتمال می‌تواند از $\tfrac12$ بیشتر باشد؟
\answer{$\P(X\ \fa{زوج})=\sum_{j\ge1}(1-p)^{2j-1}p=\frac{1-p}{2-p}<\tfrac12$؛ چون همیشه $p_X(1)>p_X(2)$، $p_X(3)>p_X(4)$، \dots}
\end{exercise}

\begin{exercise}
علی و رضا به نوبت سکه‌ای را که با احتمال $p$ شیر می‌آید پرتاب می‌کنند و اولین کسی که شیر بیاورد برنده است. علی شروع می‌کند. احتمال برد علی را بیابید. برای سکهٔ سالم چقدر است؟
\answer{علی برنده است اگر و تنها اگر زمان اولین شیر فرد باشد: $\frac{1}{2-p}$؛ برای $p=\tfrac12$ برابر $\tfrac23$.}
\end{exercise}

\begin{exercise}
هر بلیت بخت‌آزمایی با احتمال $\frac1{1000}$ برنده است. کسی هر روز یک بلیت می‌خرد. احتمال اینکه در ۱۰۰۰ روز اول هیچ‌وقت برنده نشود چقدر است؟ این عدد را با $e^{-1}$ مقایسه کنید و ارتباطش را با پواسون توضیح دهید.
\answer{$\P(X>1000)=0.999^{1000}\approx0.368\approx e^{-1}$؛ همان $\P(Y=0)$ برای $Y\sim\mathrm{Poisson}(1)$، چون تعداد بردها $\mathrm{Bin}(1000,0.001)\approx\mathrm{Poisson}(1)$ است.}
\end{exercise}

\begin{exercise}
دو نفر مستقل از هم هر کدام سکهٔ خودشان را پرتاب می‌کنند (احتمال شیر به ترتیب $p$ و $r$) تا شیر بیاورند؛ زمان‌ها $X$ و $Y$ هستند. توزیع $\min(X,Y)$، یعنی اولین لحظه‌ای که دست‌کم یکی شیر آورده، را بیابید. (راهنمایی: $\P(\min(X,Y)>k)$.)
\answer{$\P(\min>k)=(1-p)^k(1-r)^k$؛ پس $\min(X,Y)\sim\mathrm{Geom}\bigl(1-(1-p)(1-r)\bigr)$.}
\end{exercise}

\begin{exercise}
تاسی را آن‌قدر می‌ریزیم تا ۶ بیاید. به شرط اینکه ۶ در سه پرتاب اول نیامده باشد، احتمال اینکه دقیقاً در پرتاب هفتم بیاید چقدر است؟ آن را یک بار مستقیم از تعریف احتمال شرطی و یک بار با بی‌حافظگی حساب کنید.
\answer{$\P(X=7\mid X>3)=\P(X=4)=(5/6)^3\cdot\tfrac16\approx0.096$}
\end{exercise}

\begin{exercise}
نشان دهید اگر $\lambda$ عدد صحیح باشد، $\mathrm{Poisson}(\lambda)$ دو قله در $\lambda-1$ و $\lambda$ دارد. برای $\lambda=2.5$ قله کجاست؟
\answer{$p_X(\lambda)/p_X(\lambda-1)=\lambda/\lambda=1$؛ برای $\lambda=2.5$ قله در $k=2$ است.}
\end{exercise}

\begin{exercise}
تعداد غلط‌های تایپی هر صفحه $\mathrm{Poisson}(\tfrac12)$ است. احتمال اینکه یک صفحه دست‌کم ۲ غلط داشته باشد چقدر است؟
\answer{$1-e^{-1/2}(1+\tfrac12)\approx0.090$}
\end{exercise}

\begin{exercise}
یک کارخانه تراشه‌هایی می‌سازد که هر کدام مستقلاً با احتمال $0.001$ معیوب است. در محموله‌ای ۲۰۰۰ تایی، احتمال اینکه حداکثر ۲ تراشهٔ معیوب باشد را با تقریب پواسون حساب کنید.
\answer{$\lambda=2$ \ و \ $e^{-2}(1+2+2)=5e^{-2}\approx0.677$ (مقدار دقیق دوجمله‌ای هم تا سه رقم همین است).}
\end{exercise}

\begin{exercise}
اگر $X\sim\mathrm{Poisson}(\lambda)$، نشان دهید $\P(X\ \fa{زوج})=\frac{1+e^{-2\lambda}}{2}$. (راهنمایی: $e^{\lambda}+e^{-\lambda}$ را با سری تیلور بنویسید. این نتیجه را با مثال «کلید چراغ» جلسهٔ ۵ مقایسه کنید.)
\answer{$\sum_{k\ \fa{زوج}}e^{-\lambda}\frac{\lambda^k}{k!}=e^{-\lambda}\cdot\frac{e^\lambda+e^{-\lambda}}{2}$. با $p=\lambda/n$ در جواب کلید چراغ، $\frac{1+(1-2\lambda/n)^n}{2}\to\frac{1+e^{-2\lambda}}{2}$.}
\end{exercise}

\begin{exercise}\hard
اگر $X\sim\mathrm{Poisson}(\lambda)$ و $Y\sim\mathrm{Poisson}(\mu)$ و پیشامدهای مربوط به $X$ از پیشامدهای مربوط به $Y$ مستقل باشند، نشان دهید $X+Y\sim\mathrm{Poisson}(\lambda+\mu)$. (راهنمایی: $\P(X+Y=n)=\sum_{k=0}^n\P(X=k)\P(Y=n-k)$ و بسط دوجمله‌ای.) این نتیجه را با «دو بازهٔ زمانی پشت سر هم» در فرایند پواسون توضیح دهید.
\answer{$\sum_{k=0}^ne^{-\lambda}\frac{\lambda^k}{k!}e^{-\mu}\frac{\mu^{n-k}}{(n-k)!}=\frac{e^{-(\lambda+\mu)}}{n!}\sum_k\binom nk\lambda^k\mu^{n-k}=e^{-(\lambda+\mu)}\frac{(\lambda+\mu)^n}{n!}$}
\end{exercise}
